\documentclass[11pt]{amsart}
\usepackage[T1]{fontenc}
\usepackage{lmodern,microtype,amsmath,amssymb,amsthm,mathtools,enumitem}
\usepackage[margin=1.1in]{geometry}
\usepackage[hidelinks]{hyperref}
\usepackage{xurl}

\newtheorem{theorem}{Theorem}[section]
\newtheorem{lemma}[theorem]{Lemma}
\newtheorem{corollary}[theorem]{Corollary}
\newtheorem{proposition}[theorem]{Proposition}
\theoremstyle{remark}
\newtheorem{remark}[theorem]{Remark}
\numberwithin{equation}{section}

\newcommand{\E}{\mathbb E}
\newcommand{\Pp}{\mathbb P}
\newcommand{\1}{\mathbf 1}
\newcommand{\F}{\mathcal F}

\hypersetup{pdftitle={Filled-in peeling explorations exhaust every critical Boltzmann planar map},pdfauthor={Yitzchak Shmalo}}

\title[Filled-in peeling explorations exhaust critical Boltzmann maps]{Filled-in peeling explorations exhaust\\ every critical Boltzmann planar map}
\author{Yitzchak Shmalo}
\address{Einstein Institute of Mathematics, The Hebrew University of Jerusalem, Jerusalem, Israel}
\email{yitzchak.shmalo@gmail.com}
\date{5 November 2026}
\subjclass[2020]{60D05, 05C80, 60G42, 60J10}
\keywords{random planar maps, Boltzmann maps, peeling process, filled-in exploration, random walk conditioned to stay positive}

\begin{document}
\begin{abstract}
Let $\mathbf q$ be an admissible critical weight sequence for bipartite Boltzmann planar maps and consider the infinite $\mathbf q$-Boltzmann map with a boundary. Curien proved that a filled-in peeling exploration driven by any algorithm eventually discovers the whole map when $\mathbf q$ is of type $a\in[3/2,5/2]$, and asked whether this holds for every critical weight sequence. We show that it does: for every admissible critical $\mathbf q$, every finite boundary length and every peeling algorithm that never stops, the filled-in exploration exhausts the map almost surely. The type assumption entered Curien's proof through the polynomial tail of the negative jumps of the perimeter process and through its scaling limit. We use instead the exact formula for the negative jumps in terms of the positive ones, which gives the rate at which a fixed boundary edge is swallowed up to constants, and a martingale argument for $\log h^\uparrow$ of the perimeter which shows that these rates have a divergent sum along the exploration.
\end{abstract}
\maketitle

\section{Introduction}\label{sec:intro}

The peeling process, introduced by Angel \cite{Ang} for the uniform infinite planar triangulation, explores a random planar map one face at a time, and for Boltzmann maps the spatial Markov property turns the perimeter of the explored region into a Markov chain whose law does not depend on the edges chosen to be peeled. Curien's Saint-Flour notes \cite{Cur,CurLNM} develop this theory for bipartite maps with a general weight sequence $\mathbf q=(q_k)_{k\ge1}$, where a face of degree $2k$ receives weight $q_k$. In the \emph{filled-in} exploration of the infinite $\mathbf q$-Boltzmann map with a boundary, each time a peeling step separates the unexplored region into a finite and an infinite part the finite part is added to the explored region, so the explored region $\bar{\mathfrak e}_n$ has a single hole at every time $n$. The edge to be peeled at each step is chosen by an algorithm $\mathcal A$ on the boundary of that hole.

Corollary~27 of \cite{Cur} states that if $\mathbf q$ is of type $a\in[3/2,5/2]$, then $\bigcup_n\bar{\mathfrak e}_n$ is the whole map almost surely, whatever the algorithm. The notes then ask (\cite[p.~57]{Cur}):

\begin{quote}
\emph{Open question.} Is the last corollary true for any critical weight sequence?
\end{quote}

For the uniform infinite planar triangulation and quadrangulation, Curien and Le Gall \cite{CLG} show that any peeling process eventually discovers the whole map, and Curien and Marzouk \cite{CM} give a polynomial bound on the probability that the origin is not swallowed by time $n$.

In Curien's proof the type enters twice. The conditional probability that a given boundary vertex is swallowed at the next step, when the hole has half-perimeter $\ell$, is of order $\ell^{1-a}$, which uses the polynomial tail $\nu(-k)\sim p_{\mathbf q}k^{-a}$ of the step distribution of the perimeter process; and the divergence of $\sum_n(S^\uparrow_n)^{1-a}$ along the perimeter process $S^\uparrow$ follows from the convergence of $n^{-1/(a-1)}S^\uparrow_n$ to a positive limit, through Jeulin's lemma \cite{Jeu}. Weight sequences with this tail and $a\in(3/2,5/2)$ are the non-generic critical sequences with large faces of Le Gall and Miermont \cite{LGM}; the geometry of the corresponding infinite maps was studied by Budd and Curien \cite{BC}, and the case $a=2$ by Budd, Curien and Marzouk \cite{BCM}. Neither the tail asymptotics nor the scaling limit is available for a general critical weight sequence. We prove that they are not needed.

\begin{theorem}\label{thm:main}
Let $\mathbf q$ be an admissible critical weight sequence, let $\ell\ge1$, and let $\mathcal A$ be a peeling algorithm that never stops, possibly using additional randomness independent of the map. Under $\mathbb P^{(\ell)}_\infty$ the filled-in exploration $(\bar{\mathfrak e}_n)_{n\ge0}$ driven by $\mathcal A$ satisfies $\bigcup_{n\ge0}\bar{\mathfrak e}_n=\mathfrak m$ almost surely.
\end{theorem}

The algorithm in \cite{Cur} is allowed to take a cemetery value, and an exploration that stops at a finite time explores a finite region; Corollary~27 and its question concern explorations that go on forever, which is the hypothesis of Theorem~\ref{thm:main}.

The proof has three ingredients. The first is an exact formula for the negative jumps of the perimeter process in terms of the positive ones (Proposition~\ref{prop:kernel}):
\[
\nu(-k)=\frac{d(k)}{2k-1}\sum_{m\ge1}\frac{m(2m+1)d(m)\nu(m)}{m+k},\qquad d(k)=4^{-k}\binom{2k}k .
\]
In the full-perimeter normalisation this formula is due to Budd \cite{Bud}; we derive it directly from the integral representation of the disk partition functions in \cite{Cur}, which shows that it holds for every critical weight sequence without any regularity assumption. With $T(p)$ the sum on the right at $k=p$ and $R(p)=p^{-1/2}T(p)$, the formula shows that the swallowing probability at half-perimeter $p$ is at least a constant times $R(p)$ (Lemma~\ref{lem:est}). For a weight sequence of type $a$ one has $R(p)\asymp p^{1-a}$, Curien's rate. The same quantity controls the probability of a jump of relative size more than one half and the mean and variance of the truncated increments of $\log h^\uparrow(P_n)$. The second ingredient is the identity $\E_p[h^\uparrow(p)/h^\uparrow(P_1)]=1-\delta(p)$, where $\delta(p)$ is the probability that the underlying walk jumps below $1$ from $p$; it gives the cancellation between upward and downward steps that a direct estimate of the drift of $\log h^\uparrow(P_n)$ would miss. The third is a localisation argument: if $\sum_nR(P_n)$ were finite on an event of positive probability, then on that event $\log h^\uparrow(P_n)$ would converge, while the perimeter tends to infinity (Proposition~\ref{prop:sum}). The proof is completed by a supermartingale bound for the survival of a fixed boundary edge and a path argument in the dual map.

Section~\ref{sec:setup} recalls the setting. Section~\ref{sec:kernel} proves the kernel formula and the estimates that follow from it, Section~\ref{sec:perimeter} studies the perimeter process, Section~\ref{sec:proof} proves Theorem~\ref{thm:main}, and Section~\ref{sec:numerics} describes the computer checks that accompany the paper. Numbered statements of \cite{Cur} refer to the version of the notes on the author's web page; the published book is \cite{CurLNM}.

\section{Setting}\label{sec:setup}

We follow \cite{Cur}. Let $\mathbf q=(q_k)_{k\ge1}$ be a sequence of nonnegative numbers and
\[
f(x)=\sum_{k\ge1}\binom{2k-1}{k-1}q_kx^{k-1}.
\]
The sequence is admissible when the equation $f(x)=1-1/x$ has a positive solution; let $Z$ be the smallest one. It is critical when moreover $Z^2f'(Z)=1$ \cite[Proposition~5, Definition~9]{Cur}. Since $f'$ only involves $q_k$ with $k\ge2$, criticality forces $q_k>0$ for some $k\ge2$. Put $c=4Z$,
\[
d(j)=4^{-j}\binom{2j}j\quad(j\ge0),\qquad h(j)=2j\,d(j)\quad(j\ge1),\qquad h(j)=0\quad(j\le0),
\]
so that $h$ is the function $h^\uparrow$ of \cite{Cur}. Let $W^{(j)}$ be the partition function of $\mathbf q$-Boltzmann maps with a boundary of perimeter $2j$, with $W^{(0)}=1$, and define a measure $\nu$ on $\mathbb Z$ by
\[
\nu(m)=q_{m+1}c^m\quad(m\ge0),\qquad\nu(-k)=2W^{(k-1)}c^{-k}\quad(k\ge1).
\]
For admissible critical $\mathbf q$, $\nu$ is a probability measure and $h$ is harmonic for the $\nu$-random walk killed when it leaves $\{1,2,\dots\}$ \cite[Lemma~11]{Cur}:
\begin{equation}\label{eq:harm}
\sum_{j\in\mathbb Z}\nu(j)h(p+j)=h(p),\qquad p\ge1 .
\end{equation}
We use the elementary facts
\begin{equation}\label{eq:dh}
\frac1{2\sqrt j}\le d(j)\le\frac1{\sqrt j},\qquad \sqrt j\le h(j)\le2\sqrt j,\qquad \frac{h(j+1)}{h(j)}=1+\frac1{2j}\qquad(j\ge1).
\end{equation}

Let $\mathfrak m$ be a map with law $\mathbb P^{(\ell)}_\infty$, the infinite $\mathbf q$-Boltzmann map with a boundary of perimeter $2\ell$ \cite[Theorem~8]{Cur}. The filled-in exploration starts from $\bar{\mathfrak e}_0$, the boundary face. At each step $\mathcal A$ selects an edge on the boundary of the hole of $\bar{\mathfrak e}_n$, which is peeled; if the step encloses a finite region, that region is filled in. Let $\F_n$ be the $\sigma$-field generated by $\bar{\mathfrak e}_0,\dots,\bar{\mathfrak e}_n$ and the randomness used by $\mathcal A$ up to time $n$, and let $P_n$ be the half-perimeter of the hole of $\bar{\mathfrak e}_n$, so $P_0=\ell$. The boundary of the hole is a cyclic sequence of $2P_n$ edges of $\mathfrak m$, each of which has a face of $\bar{\mathfrak e}_n$ on its other side; we call these edges \emph{exposed} at time $n$. Number the positions of the cyclic sequence from the selected edge, at position $0$, to the left by $1,2,\dots,2P_n-1$; the edge at position $i$ is at position $2P_n-i$ from the right. The spatial Markov property gives the following description of one step \cite[Section~4.2]{Cur}. Conditionally on $\F_n$, with $p=P_n$:
\begin{enumerate}[label=(P\arabic*),leftmargin=*]
\item for each $k\ge1$, the peeled edge is incident to a new face of degree $2k$ with probability $\nu(k-1)h(p+k-1)/h(p)$; then $P_{n+1}=p+k-1$;
\item for each $0\le k\le p-2$, the peeled edge is glued to the edge at position $2k+1$ on its left, and the $2k$ edges strictly between them bound a finite region which is filled in, with probability $\frac12\nu(-k-1)h(p-k-1)/h(p)$; then $P_{n+1}=p-k-1$. The same holds with left replaced by right.
\end{enumerate}
These probabilities do not depend on the selected edge, and by \eqref{eq:harm} they add up to one. In particular $(P_n)$ is a Markov chain for $(\F_n)$ with transition kernel
\begin{equation}\label{eq:K}
K(p,y)=\nu(y-p)\frac{h(y)}{h(p)},\qquad p,y\ge1 .
\end{equation}
In each case the selected edge stops being exposed, and in case (P2) so do the $2k+1$ edges at positions $1,\dots,2k+1$ on the side of the gluing. An edge that has stopped being exposed has faces of $\bar{\mathfrak e}_n$ on both sides and is never exposed again.

\section{The negative jumps}\label{sec:kernel}

Put
\[
w_m=m(2m+1)d(m)\nu(m)\ \ (m\ge1),\qquad T(x)=\sum_{m\ge1}\frac{w_m}{m+x}\ \ (x>0),\qquad R(p)=\frac{T(p)}{\sqrt p}\ \ (p\ge1).
\]

\begin{proposition}\label{prop:kernel}
For every admissible critical $\mathbf q$ and every $k\ge1$,
\[
\nu(-k)=\frac{d(k)}{2k-1}\,T(k).
\]
Moreover $T(1)=1/Z$, and $0<T(x)<\infty$ for every $x>0$.
\end{proposition}

\begin{proof}
For $0\le g\le1$ let $\mathbf q_g=(g^{k-1}q_k)_{k\ge1}$. Equation (5.3) of \cite{Cur} expresses the disk partition functions as
\begin{equation}\label{eq:W}
W^{(j)}=d(j)\int_0^1\bigl(g\,c_{\mathbf q_g}\bigr)^j\,dg,\qquad j\ge0,
\end{equation}
where $c_{\mathbf q_g}=4Z_{\mathbf q_g}$. We change variables. Put $a_k=q_k\binom{2k-1}{k-1}Z^{k-1}$, so that $f(Zr)=\sum_{k\ge1}a_kr^{k-1}$, and let $J(r)=Zr\,(1-f(Zr))$ for $0\le r\le1$. The two critical equations give $\sum_{k\ge1}ka_k=f(Z)+Zf'(Z)=1$, hence
\[
J'(r)=Z\Bigl(1-\sum_{k\ge1}ka_kr^{k-1}\Bigr)=Z\sum_{k\ge2}ka_k\bigl(1-r^{k-1}\bigr),\qquad0\le r\le1 .
\]
The last series has nonnegative terms and converges uniformly on $[0,1]$, and it is positive for $r<1$ because some $a_k$ with $k\ge2$ is positive. So $J$ increases continuously from $J(0)=0$ to $J(1)=Z(1-f(Z))=1$. The weight sequence $\mathbf q_g$ satisfies $f_{\mathbf q_g}(x)=f(gx)$, so $x$ solves $f_{\mathbf q_g}(x)=1-1/x$ if and only if $y=gx$ solves $y(1-f(y))=g$. On $[0,Z]$ the function $y(1-f(y))=J(y/Z)$ is increasing, so the smallest solution is $y=ZJ^{-1}(g)$, and $g\,c_{\mathbf q_g}=4ZJ^{-1}(g)=c\,J^{-1}(g)$. No property of $f$ beyond $Z$ is used. The substitution $g=J(r)$ in \eqref{eq:W} gives
\[
W^{(j)}=c^j d(j)\int_0^1r^jJ'(r)\,dr .
\]
For $k\ge1$, by Tonelli's theorem,
\[
\int_0^1r^{k-1}J'(r)\,dr=Z\sum_{l\ge2}la_l\Bigl(\frac1k-\frac1{k+l-1}\Bigr)=\frac Zk\sum_{m\ge1}\frac{m(m+1)a_{m+1}}{m+k}.
\]
Since $\binom{2m+1}{m}=\frac{2m+1}{m+1}\binom{2m}m$ and $\nu(m)=q_{m+1}4^mZ^m$, we have $(m+1)a_{m+1}=(2m+1)d(m)\nu(m)$, so the last sum is $T(k)$. Therefore
\[
\nu(-k)=\frac2{c^k}W^{(k-1)}=\frac2c\,d(k-1)\,\frac Zk\,T(k)=\frac{d(k-1)}{2k}T(k)=\frac{d(k)}{2k-1}T(k),
\]
using $c=4Z$ and $2k\,d(k)=(2k-1)d(k-1)$. For $k=1$ the integral equals $J(1)-J(0)=1$, so $T(1)=1/Z$. Finally $T$ is nonincreasing, so $T(x)\le T(1)<\infty$ for $x\ge1$, and $T(x)\le(1+1/x)T(1)$ for $0<x<1$; and $T(x)>0$ because $\nu(m)>0$ for some $m\ge1$.
\end{proof}

In particular $\nu(-k)>0$ for every $k\ge1$. By \eqref{eq:dh}, for $k,m\ge1$,
\begin{equation}\label{eq:compar}
\tfrac14k^{-3/2}T(k)\le\nu(-k)\le k^{-3/2}T(k),\qquad m^{3/2}\nu(m)\le w_m\le3m^{3/2}\nu(m),
\end{equation}
and comparing the series term by term,
\begin{equation}\label{eq:Tmono}
T(p)\le T(k)\le\frac pkT(p)\quad(1\le k\le p),\qquad T(k)\le T(p)\quad(k\ge p).
\end{equation}
Both factors of $R(p)$ are nonincreasing, so $R(p)\le R(1)=1/Z$. Let
\[
\delta(p)=\sum_{k\ge p}\nu(-k),\qquad b(p)=\frac12\sum_{k=\lceil p/2\rceil}^{p-2}\nu(-k-1)\frac{h(p-k-1)}{h(p)},
\]
the probability that the $\nu$-walk jumps from $p$ to a point below $1$, and the probability in (P2) that a gluing on one prescribed side encloses a region of half-perimeter at least $p/2$.

\begin{lemma}\label{lem:est}
There is an absolute constant $C$ such that for every admissible critical $\mathbf q$ and every $p\ge1$:
\begin{enumerate}[label=\textup{(\alph*)},leftmargin=*]
\item $\sum_{k>p/2}\nu(-k)\le C\,R(p)$; in particular $\delta(p)\le C\,R(p)$;
\item $\sum_{1\le k\le p}k^2\nu(-k)\le C\,p^2R(p)$;
\item $\sum_{1\le m\le p}m^2\nu(m)\le C\,p^2R(p)$;
\item $\sum_{m>p/2}\sqrt{m/p}\;\nu(m)\le C\,R(p)$.
\end{enumerate}
Moreover $b(p)\ge R(p)/256$ for $p\ge16$, and $b(p)>0$ for $p\ge4$.
\end{lemma}

\begin{proof}
(a) For $k>p/2$, \eqref{eq:Tmono} gives $T(k)\le2T(p)$, so by \eqref{eq:compar} the sum is at most $2T(p)\sum_{k\ge k_0}k^{-3/2}$ with $k_0=\lfloor p/2\rfloor+1>p/2$, and $\sum_{k\ge k_0}k^{-3/2}\le k_0^{-3/2}+2k_0^{-1/2}\le3\sqrt2\,p^{-1/2}$.

(b) By \eqref{eq:compar} and \eqref{eq:Tmono}, $\sum_{k\le p}k^2\nu(-k)\le\sum_{k\le p}k^{1/2}\frac pkT(p)\le2p^{3/2}T(p)$.

(c) By \eqref{eq:compar}, $m^2\nu(m)\le m^{1/2}w_m\le p^{1/2}w_m$ for $m\le p$, and $\sum_{m\le p}w_m\le2p\sum_{m\le p}w_m/(m+p)\le2pT(p)$.

(d) For $m>p/2$ we have $m+p<3m$, so $w_m/(m+p)\ge w_m/(3m)\ge\frac13m^{1/2}\nu(m)$ and $\sum_{m>p/2}m^{1/2}\nu(m)\le3T(p)$.

For the lower bound on $b$, let $p\ge16$ and restrict the sum to $j=k+1$ with $\lceil p/2\rceil+1\le j\le\lfloor3p/4\rfloor$. There are at least $p/4-2\ge p/8$ such $j$, each satisfies $k\le p-2$, $j\le p$ and $p-j\ge p/4$. By \eqref{eq:compar} and \eqref{eq:Tmono}, $\nu(-j)\ge\frac14p^{-3/2}T(p)$, and by \eqref{eq:dh}, $h(p-j)/h(p)\ge\frac{\sqrt{p/4}}{2\sqrt p}=\frac14$. Hence $b(p)\ge\frac12\cdot\frac p8\cdot\frac{T(p)}{4p^{3/2}}\cdot\frac14=R(p)/256$. For $p\ge4$ the sum defining $b(p)$ is not empty and all its terms are positive.
\end{proof}

\section{The perimeter process}\label{sec:perimeter}

For $p\ge1$ let $S_p=\{y\ge1:|y-p|\le p/2\}$ and $B_p=\{y\ge1:|y-p|>p/2\}$, and write $\E_p$ for the law of the chain \eqref{eq:K} started at $p$.

\begin{lemma}\label{lem:big}
With an absolute constant $C$, for every $p\ge1$,
\[
K(p,B_p)\le C\,R(p)\qquad\text{and}\qquad\sum_{y\in B_p}\nu(y-p)\le C\,R(p).
\]
\end{lemma}

\begin{proof}
For $y=p-k\in B_p$ we have $k>p/2$ and $h(y)\le h(p)$, so both sums over these $y$ are at most $\sum_{k>p/2}\nu(-k)$, which is bounded in Lemma~\ref{lem:est}(a). For $y=p+m\in B_p$ we have $m>p/2$, hence $p+m<3m$, $h(y)/h(p)\le2\sqrt{(p+m)/p}\le2\sqrt3\sqrt{m/p}$ by \eqref{eq:dh}, and $1\le\sqrt2\sqrt{m/p}$; Lemma~\ref{lem:est}(d) bounds both sums over these $y$.
\end{proof}

\begin{proposition}\label{prop:escape}
Almost surely $P_n\to\infty$.
\end{proposition}

\begin{proof}
By \eqref{eq:K}, $\E_p[1/h(P_1)]=\sum_{y\ge1}\nu(y-p)/h(p)=(1-\delta(p))/h(p)$, so $1/h(P_n)$ is a positive supermartingale and
\[
\E\sum_{n\ge0}\frac{\delta(P_n)}{h(P_n)}\le\frac1{h(\ell)}<\infty .
\]
Hence the sum is finite almost surely. Since $\delta(p)>0$ for every $p$ by Proposition~\ref{prop:kernel}, every state is visited finitely often almost surely, and $P_n\to\infty$.
\end{proof}

Let $U=\log h$ on $\{1,2,\dots\}$, and for a transition from $p$ to $y$ put
\[
L(p,y)=\bigl(U(y)-U(p)\bigr)\1_{S_p}(y).
\]

\begin{lemma}\label{lem:L}
For every $p\ge1$ and $y\in S_p$ we have $|U(y)-U(p)|\le|y-p|/p\le1/2$. With an absolute constant $C$,
\[
\E_p\bigl[L(p,P_1)^2\bigr]\le C\,R(p)\qquad\text{and}\qquad\bigl|\E_p[L(p,P_1)]\bigr|\le C\,R(p).
\]
\end{lemma}

\begin{proof}
By \eqref{eq:dh}, $|U(y)-U(p)|$ is a sum of $|y-p|$ terms $\log(1+1/(2j))$ with $j\ge\min(y,p)\ge p/2$, so it is at most $|y-p|/p$. In particular $h(y)/h(p)\le e^{1/2}<2$ on $S_p$, and Lemma~\ref{lem:est}(b),(c) give
\[
\E_p[L^2]\le\sum_{y\in S_p}\nu(y-p)\frac{h(y)}{h(p)}\frac{(y-p)^2}{p^2}\le\frac2{p^2}\Bigl(\sum_{k\le p}k^2\nu(-k)+\sum_{m\le p}m^2\nu(m)\Bigr)\le C\,R(p).
\]
For the mean, put $H=h(P_1)/h(p)$. Then
\[
\E_p\bigl[1-H^{-1}\bigr]=1-\sum_{y\ge1}\nu(y-p)=\delta(p),
\]
and
\[
\E_p\bigl[|1-H^{-1}|\,\1_{B_p}(P_1)\bigr]=\sum_{y\in B_p}\nu(y-p)\Bigl|\frac{h(y)}{h(p)}-1\Bigr|\le K(p,B_p)+\sum_{y\in B_p}\nu(y-p).
\]
By Lemma~\ref{lem:est}(a) and Lemma~\ref{lem:big}, $|\E_p[(1-H^{-1})\1_{S_p}(P_1)]|\le C\,R(p)$. On $S_p$ write $u=U(P_1)-U(p)=\log H$, so $|u|\le1/2$ and $1-H^{-1}=1-e^{-u}$. Taylor's formula gives $|u-(1-e^{-u})|\le\frac12u^2e^{1/2}\le u^2$, so
\[
\bigl|\E_p[L]\bigr|\le\bigl|\E_p[(1-H^{-1})\1_{S_p}(P_1)]\bigr|+\E_p[L^2]\le C\,R(p).\qedhere
\]
\end{proof}

The identity for $\E_p[1-H^{-1}]$ is what makes the mean small: the upward and downward steps of $U(P_n)$ are each of a larger order than $R(p)$ in general, and only their difference is controlled.

\begin{proposition}\label{prop:sum}
Almost surely $\sum_{n\ge0}R(P_n)=\infty$.
\end{proposition}

\begin{proof}
Let $E=\{\sum_nR(P_n)<\infty\}$. Put $L_n=L(P_n,P_{n+1})$, $\mu_n=\E[L_n\mid\F_n]$, $D_n=L_n-\mu_n$ and $A_n=\sum_{i<n}R(P_i)$. By Lemma~\ref{lem:L}, $|D_n|\le1$, $|\mu_n|\le C\,R(P_n)$ and $\E[D_n^2\mid\F_n]\le C\,R(P_n)$. For $K>0$ let $\tau_K=\inf\{n:A_n>K\}$; since $A_n$ is $\F_{n-1}$-measurable, $\{n<\tau_K\}\in\F_n$. As $R\le R(1)$,
\[
\sum_{n<\tau_K}R(P_n)\le K+R(1).
\]
The process $\sum_{n<N}\1\{n<\tau_K\}D_n$ is a martingale with
\[
\E\Bigl[\Bigl(\sum_{n<N}\1\{n<\tau_K\}D_n\Bigr)^2\Bigr]=\E\sum_{n<N}\1\{n<\tau_K\}\E[D_n^2\mid\F_n]\le C\bigl(K+R(1)\bigr),
\]
so it converges almost surely as $N\to\infty$. Similarly, by Lemma~\ref{lem:big},
\[
\E\sum_{n<\tau_K}\1\{P_{n+1}\in B_{P_n}\}=\E\sum_{n<\tau_K}K(P_n,B_{P_n})\le C\bigl(K+R(1)\bigr),
\]
so only finitely many $n<\tau_K$ have $P_{n+1}\in B_{P_n}$, almost surely. On $E$ we have $\tau_K=\infty$ for every integer $K>A_\infty$. Taking the union over integers $K$, almost surely on $E$: $\sum_nD_n$ converges, $\sum_n|\mu_n|\le C\sum_nR(P_n)<\infty$, and $U(P_{n+1})-U(P_n)\ne L_n$ for only finitely many $n$. Each of these finitely many increments is finite, so $U(P_N)=U(\ell)+\sum_{n<N}(U(P_{n+1})-U(P_n))$ converges to a finite limit on $E$. But $U(P_N)=\log h(P_N)\to\infty$ by Proposition~\ref{prop:escape}. Hence $\Pp(E)=0$.
\end{proof}

\begin{remark}
When $\nu(-k)\sim p_{\mathbf q}k^{-a}$ with $a\in(3/2,5/2]$, \eqref{eq:compar} gives $T(p)\asymp p^{3/2-a}$ and $R(p)\asymp p^{1-a}$, and Proposition~\ref{prop:sum} becomes the divergence of $\sum_n(S^\uparrow_n)^{1-a}$, which \cite{Cur} obtains from the scaling limit of $S^\uparrow$ and Jeulin's lemma.
\end{remark}

\section{Proof of Theorem~\ref{thm:main}}\label{sec:proof}

\begin{lemma}\label{lem:swallow}
If an edge $\varepsilon$ is exposed at time $n$ and $P_n\ge4$, then $\Pp(\varepsilon\text{ is not exposed at time }n+1\mid\F_n)\ge b(P_n)$.
\end{lemma}

\begin{proof}
Let $p=P_n$. If $\varepsilon$ is the selected edge, it stops being exposed with probability one. Otherwise it is at a position $i\in\{1,\dots,2p-1\}$ on the left, that is at position $2p-i$ on the right, and one of $i,2p-i$ is at most $p$; let it be the position $i'$ on the side $\sigma$. By (P2), for $\lceil p/2\rceil\le k\le p-2$ the selected edge is glued to the edge at position $2k+1$ on the side $\sigma$ with probability $\frac12\nu(-k-1)h(p-k-1)/h(p)$, and then the edges at positions $1,\dots,2k+1\ge p+1>i'$ on that side stop being exposed. These events are disjoint and their probabilities add up to $b(p)$.
\end{proof}

\begin{lemma}\label{lem:hazard}
Let $(\mathcal G_n)_{n\ge0}$ be a filtration and $\tau$ a stopping time with $\tau\ge1$. Put $r_n=\Pp(\tau=n+1\mid\mathcal G_n)$. Then $\sum_{n\ge0}r_n<\infty$ almost surely on $\{\tau=\infty\}$.
\end{lemma}

\begin{proof}
Let $Y_m=\1\{\tau>m\}\exp(\sum_{n<m}r_n)$, so $Y_0=1$. The sum in the exponent is $\mathcal G_{m-1}$-measurable, and $\Pp(\tau>m+1\mid\mathcal G_m)=1-r_m$ on $\{\tau>m\}$. Hence $\E[Y_{m+1}\mid\mathcal G_m]=Y_me^{r_m}(1-r_m)\le Y_m$, and $(Y_m)$ is a nonnegative supermartingale. It converges almost surely to a finite limit, and on $\{\tau=\infty\}$ we have $Y_m=\exp(\sum_{n<m}r_n)$.
\end{proof}

\begin{proof}[Proof of Theorem~\ref{thm:main}]
By Propositions~\ref{prop:escape} and~\ref{prop:sum} and Lemma~\ref{lem:est}, almost surely $P_n\ge16$ for all large $n$ and
\begin{equation}\label{eq:bsum}
\sum_{n\ge0}b(P_n)\1\{P_n\ge4\}=\infty .
\end{equation}
For $n\ge0$ and $i\ge0$, on the event $\{i<2P_n\}$ let $\varepsilon_{n,i}$ be the exposed edge at position $i$ at time $n$, and let $\tau_{n,i}$ be the first time after $n$ at which $\varepsilon_{n,i}$ is not exposed ($\tau_{n,i}=\infty$ if there is none). Put $\tau_{n,i}=n+1$ on $\{i\ge2P_n\}$. Then $\tau_{n,i}$ is a stopping time for $(\F_m)$, and since an edge that stops being exposed is never exposed again, Lemma~\ref{lem:swallow} gives $\Pp(\tau_{n,i}=m+1\mid\F_m)\ge b(P_m)\1\{P_m\ge4\}$ on $\{\tau_{n,i}>m\ge n\}$. By Lemma~\ref{lem:hazard} and \eqref{eq:bsum}, $\Pp(\tau_{n,i}=\infty)=0$. Every edge that is ever exposed is one of the countably many $\varepsilon_{n,i}$, so almost surely every edge that is ever exposed stops being exposed at a finite time.

On this event, let $f$ be a face of $\mathfrak m$ and let $f_0,f_1,\dots,f_r=f$ be a path in the dual map with $f_0$ the boundary face, consecutive faces sharing an edge. If some $f_j$ is not in $\bigcup_n\bar{\mathfrak e}_n$, take the smallest such $j$; then $j\ge1$ and $f_{j-1}\in\bar{\mathfrak e}_s$ for some $s$. The edge shared by $f_{j-1}$ and $f_j$ is an edge of $\bar{\mathfrak e}_s$ whose other side is not in $\bar{\mathfrak e}_s$, so it is exposed at time $s$. It stops being exposed at some finite time $m$, and then it has faces of $\bar{\mathfrak e}_m$ on both sides, so $f_j\in\bar{\mathfrak e}_m$, a contradiction. Thus every face of $\mathfrak m$ belongs to $\bigcup_n\bar{\mathfrak e}_n$, and so does every edge and vertex, each being incident to a face.
\end{proof}

\section{Computations}\label{sec:numerics}

The repository \cite{code} contains three Python programs. The first, \texttt{kernel\_check.py}, tests Proposition~\ref{prop:kernel} on the critical weight sequence
\[
q_1=0,\qquad q_k=\frac{C_*\,k^{-4}Z_*^{1-k}}{\binom{2k-1}{k-1}}\quad(k\ge2),\qquad C_*=\frac1{\zeta(3)-1},\qquad Z_*=\frac{\zeta(3)-1}{\zeta(3)-\zeta(4)},
\]
for which $f(Z_*r)=C_*\sum_{k\ge2}k^{-4}r^{k-1}$, so that $f(Z_*)=C_*(\zeta(4)-1)=1-1/Z_*$ and $Z_*^2f'(Z_*)=Z_*C_*(\zeta(3)-\zeta(4))=1$; here $f$ has radius of convergence exactly $Z_*$. In this case $J'(r)=Z_*C_*(\zeta(3)-\mathrm{Li}_3(r)/r)$, and $\int_0^1r^jJ'(r)\,dr=Z_*C_*\bigl(\zeta(2)/j^2-\zeta(3)/(j(j+1))-H_j/j^3\bigr)$ for $j\ge1$, with $H_j$ the harmonic number. The program checks the two critical equations, compares $\nu(-k)$ from this closed form with $d(k)T(k)/(2k-1)$, checks that $\nu$ has total mass one, and checks \eqref{eq:harm} for $p\le6$. The relative differences are below $2\cdot10^{-40}$ in $40$-digit arithmetic for $k\le50$ and at the level of double-precision rounding for $k=200$ and $k=1000$; the total mass and the harmonicity agree to the same accuracy.

\sloppy
The second program, \texttt{universal\_checks.py}, takes positive laws $\nu(m)$, $m\ge1$, with $\sum_m\nu(m)h(m+1)=1$, concentrated on single atoms $m\in\{1,4,64,4096,10^6\}$ or spread over these five atoms, defines $\nu(-k)$ by Proposition~\ref{prop:kernel}, and computes the rows of \eqref{eq:K} at $p=1,2,4,8,16$ and $p=4^j$, $3\le j\le7$. All $70$ rows sum to one within $6\cdot10^{-16}$, so the kernel formula reproduces \eqref{eq:harm} for these laws. Over these rows the ratios to $R(p)$ of $K(p,B_p)$, $|\E_p L|$, $\E_pL^2$ and $\delta(p)$ are at most $2.66$, $0.51$, $0.30$ and $1.00$, and $b(p)/R(p)\ge0.044$ for $p\ge8$.

A third program, \texttt{experiment.py}, verifies the closed form of $\int_0^1r^jJ'(r)\,dr$ against numerical quadrature for $j\le10$ (to $2\cdot10^{-56}$), shows that $j^2\int_0^1r^jJ'(r)\,dr$ approaches $2D$ with $D=1+Z_*^3f''(Z_*)/2$ (the ratio is $0.998$ at $j=10^4$), and sums truncated rows of \eqref{eq:K} for this sequence. The repository also contains a Lean file with six elementary identities and inequalities checked in Mathlib: three about the second-order expansion of $y(1-f(y))$ at $y=Z$ under the critical equations, and three used in Section~\ref{sec:kernel}, namely $k(m+p)\le p(m+k)$ for $m\ge0$ and $k\le p$ (behind \eqref{eq:Tmono}), $m+p\le3m$ for $p\le2m$ (behind Lemma~\ref{lem:est}(d)), and the passage from $2k\,d(k)=(2k-1)d(k-1)$ to the final step of the proof of Proposition~\ref{prop:kernel}.

\section*{Acknowledgements and funding}
The research presented in this paper was supported by the European Research
Council (ERC) under the European Union's Horizon Europe research and
innovation programme (grant agreement No.~101041711), by the Simons
Foundation as part of the Collaboration on the Mathematical and Scientific
Foundations of Deep Learning, by Heights Labs, by Convex Nexus Capital, by
the Israel Science Foundation (grant number 2258/19), and by the Israel
Science Foundation (ISF Grant 4101/25). Computational resources were
provided in part by Google Cloud through its education credits program
(credit 530550103).

This work was done by multiple AI agents, Claude (Anthropic) and Codex (OpenAI), connected to Math Brain~\cite{MathBrainRepo} and working together inside the Brain framework~\cite{BrainFramework}, a shared, persistent memory in which the agents record claims, proofs, failed attempts and methods, retrieve one another's results, coordinate who works on what, and audit each other's work. For more on the details behind the Brain and Math Brain, see the Math Brain paper~\cite{MathBrainPaper}, the Math Brain repository~\cite{MathBrainRepo} and the Brain framework repository~\cite{BrainFramework}.

\section*{Declaration of competing interest}
The author holds equity in Heights Labs and in Convex Nexus Capital and
has a management role at Convex Nexus Capital. These entities are included
in the funding declaration above.

\section*{Declaration of generative AI and AI-assisted technologies}
The paper comes out of Math Brain, an ongoing project on the use of AI systems for research mathematics. In October 2026 OpenAI GPT models, running as research tasks of that project, proposed Theorem~\ref{thm:main} and its proof through the kernel formula, the estimates of Lemma~\ref{lem:est}, the localisation argument of Proposition~\ref{prop:sum} and the survival bound of Lemma~\ref{lem:hazard}, and wrote the Lean file and the program \texttt{universal\_checks.py}. In November 2026 Claude Opus~5.5 (Anthropic), working in Claude Code, checked the proof, wrote \texttt{kernel\_check.py}, and wrote the paper. The author directed the work and is responsible for the content of the paper and for decisions concerning its dissemination.

\section*{Data availability}
No data were used. The code is available at\\ \url{https://github.com/yspennstate/critical-peeling-exhaustion}.

The paper is part of the Math Brain project, whose public repository is \url{https://github.com/yspennstate/math-brain-public}~\cite{MathBrainRepo}.

\begin{thebibliography}{10}

\bibitem{Ang} O.~Angel, Growth and percolation on the uniform infinite planar triangulation, \emph{Geom. Funct. Anal.} \textbf{13} (2003), 935--974.

\bibitem{Bud}
T.~Budd, The peeling process of infinite Boltzmann planar maps, \emph{Electron. J. Combin.} \textbf{23}(1) (2016), Paper~1.28. Also arXiv:1506.01590.

\bibitem{BC}
T.~Budd and N.~Curien, Geometry of infinite planar maps with high degrees, \emph{Electron. J. Probab.} \textbf{22} (2017), 1--37.

\bibitem{BCM}
T.~Budd, N.~Curien and C.~Marzouk, Infinite random planar maps related to Cauchy processes, \emph{J. \'Ec. polytech. Math.} \textbf{5} (2018), 749--791.

\bibitem{Cur}
N.~Curien, \emph{Peeling random planar maps}, lecture notes of the \'Ecole d'\'Et\'e de Probabilit\'es de Saint-Flour XLIX (2019), preliminary version, \url{https://www.imo.universite-paris-saclay.fr/~nicolas.curien/}.

\bibitem{CurLNM}
N.~Curien, \emph{Peeling Random Planar Maps}, \'Ecole d'\'Et\'e de Probabilit\'es de Saint-Flour XLIX -- 2019, Lecture Notes in Mathematics 2335, Springer, Cham, 2023.

\bibitem{CLG} N.~Curien and J.-F.~Le Gall, Scaling limits for the peeling process on random maps, \emph{Ann. Inst. Henri Poincar\'e Probab. Stat.} \textbf{53} (2017), 322--357.

\bibitem{CM} N.~Curien and C.~Marzouk, How fast planar maps get swallowed by a peeling process, \emph{Electron. Commun. Probab.} \textbf{23} (2018), Paper No.~18.

\bibitem{Jeu}
T.~Jeulin, \emph{Semi-martingales et grossissement d'une filtration}, Lecture Notes in Mathematics 833, Springer, Berlin, 1980.

\bibitem{LGM}
J.-F.~Le Gall and G.~Miermont, Scaling limits of random planar maps with large faces, \emph{Ann. Probab.} \textbf{39} (2011), 1--69.

\bibitem{code}
Y.~Shmalo, Computations for ``Filled-in peeling explorations exhaust every critical Boltzmann planar map'', \url{https://github.com/yspennstate/critical-peeling-exhaustion}, 2026.

\bibitem{MathBrainRepo} Y.~Shmalo, \emph{Math Brain: public snapshot, code and tests}, 2026, \url{https://github.com/yspennstate/math-brain-public}.

\bibitem{MathBrainPaper} Y.~Shmalo, \emph{Math Brain: shared, auditable memory for multi-agent mathematical research}, October 2026, \url{https://github.com/yspennstate/math-brain-public/tree/main/papers/math-brain-2026-10-24}.

\bibitem{BrainFramework} Y.~Shmalo, \emph{Brain Framework}, public software repository, 2026, \url{https://github.com/yspennstate/brain-framework-public}.

\end{thebibliography}

\end{document}
